%!TEX program = xelatex
%!TEX root = geometria_analitica.tex
%%Usar makeindex -s indexstyle.ist arquivo.idx no terminal para gerar o índice remissivo agrupado por inicial
%%Após executar pdflatex arquivo

\chapter{Vetores} % (fold)
\label{cha:vetores}

% \section{Introdução} % (fold)
% \label{sec:introducao}
%
% A Geometria Analítica é um método para tratar com problemas da Geometria. No caso do plano, tal método consiste em associar os pontos à pares ordenados de números reais, permitindo assim representar curvas do plano por equações em duas variáveis, de modo que a cada curva está associado uma equação, bem definida, do tipo $f(x,y) = 0$ e a cada equação deste tipo associa-se uma curva bem definida no plano. Assim, as propriedades geométricas das curvas podem ser determinadas à partir de informações algébricas e analíticas da equação $f(x, y) = 0$.
%
% Por exemplo, a equação $5x^2 - 4xy + 8y^2 - 36 = 0$ representa uma curva em $\real^2$. Nessa forma, não podemos dizer exatamente qual curva esta equação representa. Mas utilizando-se de ferramentas da Álgebra Linear, podemos reescrever tal equação como
% \[
%   \dfrac{x'^2}{9} + \dfrac{y'^2}{4} = 1
% \]
% que é uma equação de uma elipse.

\section{Plano Cartesiano} % (fold)
\label{sec:plano_cartesiano}

Vamos considerar o plano definido pelo par de retas perpendiculares $x$ e $y$ como abaixo:
\begin{figure}[!h]
    \centering
    \caption{Plano cartesiano}
    \includegraphics{plano-cartesiano.pdf}
\end{figure}


Seja $P$ um ponto qualquer do plano. Por $P$ podemos traçar uma única reta $r$ paralela à reta $x$ e uma única reta $s$ paralela à reta $y$. Estas retas
interceptam $x$ e $y$ em pontos $P_x$ e $P_y$, respectivamente. Seja $a$ o número real correspondente a $P_x$ e $b$ o número real correspondente a $P_y$. 
\begin{figure}[!h]
    \centering
    \caption{Coordenadas do ponto $P$}
    \includegraphics{coordenadas-plano-cartesiano.pdf}
\end{figure}

Estes dois números determinam o ponto $P$, isto é, conhecendo $a$ e $b$ podemos traças as retas $r$ e $s$. O ponto $P$ é a interseção destas retas. Os números 
$a$ e $b$ são chamados \textbf{abscissa} e \textbf{ordenada} do ponto $P$, respectivamente. Eles constituem as \textbf{coordenadas} de $P$.
Denotaremos assim o ponto $P$ por
\[
    P(a,b).
\]

Sejam $P(x_1,y_1)$ e $Q(x_2,y_2)$ pontos do plano. A partir de $P$ e $Q$ podemos construir o triângulo retângulo $PQS$ como abaixo:
\begin{figure}[!h]
    \centering
    \caption{Distância entre dois pontos}
    \includegraphics{distancia-pontos-plano-cartesiano.pdf}
\end{figure}

\newpage

Assim sua hipotenusa é
\[
    \sqrt{(x_1 - x_2)^2 + (y_1 - y_2)^2}.
\]
Este número é chamado \textbf{distância} de $P$ a $Q$ e é indicado por $d(P,Q)$, isto é, por definição
\[
    d(P,Q) = \sqrt{(x_1 - x_2)^2 + (y_1 - y_2)^2}.
\]

Dizemos que dois pontos $P(x_1,y_1)$ e $Q(x_2,y_2)$ são iguais se, e somente se, $x_1 = x_2$ e $y_1 = y_2$.

% section plano_cartesiano (end)

\section{Vetores em $\real^2$} % (fold)
\label{sec:vetores_plano}

Vetores surgiram com um meio de caracterizar grandezas físicas que precisam de mais do que sua magnitude para serem
completamente determinadas. Por exemplo, grandezas com massa, área e comprimento são caracterizadas simplesmente por um
número, tais grandezas são chamadas de \textit{grandezas escalares}. No entanto, grandezas com velocidade, força e impulso só são
completamente determinadas se além da magnitude, conhecemos sua direção e sentido. Tais grandezas são chamadas de
\textit{grandezas vetorias}. Um exemplo de uma grandeza vetorias seria uma força de 4 N atuando sobre um corpo, de modo ascendente
na direção que uma um ângulo de $60\degree$ com a horizontal.

Sabemos que a cada par ordenado $(\alpha, \beta)$ corresponde um ponto do plano. Assim se $(\alpha, \beta) \ne (0,0)$ podemos associar uma seta com origem em $O(0,0)$ e extremidade em $(\alpha,\beta)$ como na figura abaixo:

\begin{figure}[!h]
    \centering
    \caption{Seta representando o ponto $(\alpha, \beta)$.}
    \includegraphics{vetores-plano-cartesiano.pdf}
\end{figure}

Assim o par ordenado $(\alpha, \beta)$ pode ser representado graficamente por um ponto ou por uma seta. Quando representamos $(\alpha, \beta)$ por uma seta,
podemos associar a este par uma \textbf{direção}, um \textbf{sentido} e um \textbf{módulo} ou \textbf{norma}. A direção do vetor determinado pelo ponto
$(\alpha, \beta)$ é a direção de uma reta passando por essa seta e o sentido é o sentido para o qual essa seta aponta. O \textbf{módulo} ou \textbf{norma} será
definido a seguir.

Em geral, um objeto ao qual podemos associar os conceitos de direção, sentido e módulo é chamado um \textbf{vetor}. Denotaremos um vetor de extremidade em um ponto $(\alpha, \beta)$ por
\[
    \vec{u} = (\alpha, \beta).
\]

Esta notação significa que $\vec{u}$ é a seta com início no ponto $O(0,0)$ e final no ponto $(\alpha, \beta)$.
\begin{figure}[!h]
    \centering
    \caption{Vetor $ \vec{u} = (\alpha, \beta)$.}
    \includegraphics{vetor-coordenadas-plano-cartesiano.pdf}
\end{figure}

% Assim escrevemos $\vec{u} = \vec{OP}$. Para simplificar a notação vamos identificar o vetor $\vec{u} = \vec{OP}$ com as coordenadas de sua extremidade e daí escrevemos
% \[
%   \vec{u} = (x_1,y_1).
% \]
% As coordenadas $(x_1,y_1)$ são chamadas de \textbf{componentes} do vetor $\vec{u}$. Com essa representação, o vetor nulo $\vec{0}$ é escrito como
% \[
%   \vec{0} = (0,0).
% \]

Por exemplo, $\vec{u} = (3,2)$ e $\vec{v} = (-1,-2)$ são representados graficamente por:
\begin{figure}[!h]
    \centering
    \caption{Vetores $\vec{u} = (3,2)$ e $\vec{v} = (-1,-2)$.}
    \includegraphics{exemplo-vetores-plano-cartesiano.pdf}
\end{figure}

O vetor que tem início e fim no ponto $(0,0)$ é chamado de \textbf{vetor nulo} e será denotado por
\[
    \vec{O} = (0,0).
\]

Podemos também representar um vetor por uma seta que não parte da origem. Por exemplo, os pontos $A(x_1,y_1)$ e $B(x_2,y_2)$ determinam o vetor
\[
    \vec{AB} = (x_2 - x_1, y_2 - y_1).
\]
\begin{figure}[!h]
    \centering
    \caption{Vetor $\vec{AB}$}
    \includegraphics{vetor-AB-plano-cartesiano.pdf}
\end{figure}

\newpage

Assim a seta que representa o vetor $\vec{AB}$ e a seta que representa o vetor $\vec{u}$, que tem início na origem e extremidade no ponto $(x_2 - x_1, y_2 - y_1)$, têm o mesmo módulo, direção e sentido que $\vec{AB}$. Assim vamos representar o vetor $\vec{AB}$ pela forma que for mais conveniente.

\begin{definicao}
    Seja $\vec{u} = (x, y)$ um vetor. A \textbf{norma} ou \textbf{norma} de $\vec{u}$, denotada por $\norm{\vec{u}}$, é dada por
    \[
        \norm{\vec{u}} = \sqrt{x^2 + y^2}.
    \]
\end{definicao}

Um vetor $\vec{u}$ tal que
\[
    \norm{\vec{u}} = 1
\]
é chamado de \textbf{vetor unitário}.
\begin{figure}[!h]
    \centering
    \caption{Vetores unitários: $\norm{\vec{u}} = 1$}
    \includegraphics{vetores-unitarios-plano-cartesiano.pdf}
\end{figure}

\begin{definicao}
    Dois vetores $\vec{u}$ e $\vec{v}$ são \textbf{paralelos} (ou \textbf{colineares}) se a restas que contém estes vetores são paralelas(podendo ser coincidentes).Caso $\vec{u}$ e $\vec{v}$ sejam paralelos, escreveremos $\vec{u}\varparallel\vec{v}$. Neste caso, $\vec{u}$ e $\vec{v}$ têm a mesma direção, mas podem ter sentidos contrários.
\end{definicao}
\begin{figure}[!h]
    \centering
    \caption{Vetores paralelos}
    \includegraphics{vetores-paralelos-plano-cartesiano.pdf}
\end{figure}

\subsection{Operações com Vetores} % (fold)
\label{sub:operacoes_com_vetores}

\begin{definicao}
  Sejam $\vec{u} = (x_1, y_1)$ e $\vec{v} = (x_2, y_2)$ vetores. Então:
  \[
      \vec{u} + \vec{v} = (x_1 + x_2, y_1 + y_2).
  \]
\end{definicao}

\begin{figure}[!h]
  \centering
  \caption{Soma de vetores}
  \includegraphics{soma-vetores-plano-cartesiano.pdf}
\end{figure}

\begin{propriedades}\label{propriedades_soma_vetores}
  Sejam $\vec{u}$, $\vec{v}$ e $\vec{w}$ vetores quaisquer. As seguintes propriedades são verdadeiras:
  \begin{enumerate}[label=({\roman*})]
    \item $(\vec{u} + \vec{v}) + \vec{w} = \vec{u} + (\vec{v} + \vec{w})$ [Associatividade]
    \item $\vec{u} + \vec{v} = \vec{v} + \vec{u}$ [Comutatividade]
    \item Existe um único vetor que somado a $\vec{u}$ resulta no próprio vetor $\vec{u}$:
    \[
      \vec{u} + \vec{0} = \vec{0} + \vec{u} = \vec{u}.
    \]
    Tal vetor é o vetor nulo. [Elemento Neutro]
    \item\label{vetor_oposto} Para cada $\vec{u}$, existe um único vetor $\vec{v}$ tal que
    \[
      \vec{u} + \vec{v} = \vec{0} = \vec{v} + \vec{u}.
    \]
    O vetor $\vec{v}$ é o vetor oposto de $\vec{u}$. [Elemento Oposto]
  \end{enumerate}
\end{propriedades}
\begin{prova}
    Sejam $\vec{u} = (x_1,y_1)$, $\vec{v} = (x_2,y_2)$ e $\vec{w} = (x_3,y_3)$. Temos
    \begin{enumerate}[label=({\roman*})]
        \item $(\vec{u} + \vec{v}) + \vec{w} = (x_1+x_2,y_1+y_2) + \vec{w} = ((x_1+x_2)+x_3,(y_1+y_2)+y_3) = (x_1+(x_2+x_3),y_1+(y_2+y_3)) = \vec{u} + (x_2+x_3,y_2+y_3) = \vec{u} + (\vec{v} + \vec{w})$.
        \item $\vec{u} + \vec{v} = (x_1+x_2,y_1+y_2) = (x_2+x_1,y_2+y_1) = \vec{v} + \vec{u}$.
        \item Usando o item (A2) basta verificar que $\vec{u} + \vec{0} = \vec{u}$. De fato
        \[
              \vec{u} + \vec{0} = (x_1,y_1) + (0,0) = (x_1,y_1) = \vec{u}.
        \]
        \item Novamente, pelo item (A2) basta verificar que $\vec{u} + \vec{v} = \vec{0}$. De fato, dado $\vec{u} = (x_1,y_1)$, tomando $\vec{v} = (-x_1,-y_1)$ temos
        \[
            \vec{u} + \vec{v} = (x_1,y_1) + (-x_1,-y_1) = (0,0) = \vec{0}.
        \]
    \end{enumerate}
    Portanto as propriedades são verdadeiras.
\end{prova}

\begin{definicao}
    O vetor $\vec{v}$ do item \ref{vetor_oposto} da Definição \ref{propriedades_soma_vetores} é chamado de \textbf{vetor oposto}.
\end{definicao}


\begin{definicao}
  Sejam $\alpha \in \real$, $\vec{u} = (x_1, y_1)$ e $\vec{v} = (x_2, y_2)$ vetores. Então:
  \[
      \alpha \vec{u} = (\alpha x_1, \alpha y_1).
  \]
\end{definicao}

\begin{figure}[!h]
  \centering
  \caption{Multiplicação por escalar}
  \includegraphics{multiplicacao-escalar-vetores-plano-cartesiano.pdf}
\end{figure}

\begin{observacao}
    \begin{enumerate}
        \item Sejam $\vec{u}$ um vetor não-nulo e $\lambda \ne 0$ um escalar. Então $\vec{u} \varparallel\lambda\vec{u}$. Além disso, se $\lambda > 0$, então $\vec{u}$ e $\lambda\vec{u}$ são de mesmo sentido, e se $\lambda < 0$ então $\vec{u}$ e $\lambda\vec{u}$ são de sentidos contrários.
        \item Quando $\lambda = -1$, o vetor $-\vec{u} = (-x_1,-y_1)$  é exatamente o vetor oposto de $\vec{u}$. Note que $-\vec{u}$ tem a mesma direção que $\vec{u}$, mas sentido contrário.
    \end{enumerate}

\end{observacao}

\begin{propriedades}
  Quaisquer que sejam os números reais $\alpha$ e $\beta$ e quaisquer que sejam os vetores $\vec{u}$ e $\vec{v}$ valem as igualdades:
  \begin{enumerate}[label=({\roman*})]
    \item $\alpha(\vec{u} + \vec{v}) = \alpha\vec{u} + \alpha\vec{v}$;
    \item $(\alpha + \beta)\vec{u} = \alpha\vec{u} + \beta\vec{u}$;
    \item $1\vec{u} = \vec{u}$;
    \item $\alpha(\beta\vec{u}) = (\alpha\beta)\vec{u} = \beta(\alpha\vec{u})$.
  \end{enumerate}
\end{propriedades}
\begin{prova}
    Sejam $\vec{u} = (x_1,y_1)$ e $\vec{v} = (x_2,y_2)$. Temos:
    \begin{enumerate}[label=({\roman*})]
        \item $\alpha(\vec{u} + \vec{v}) = \alpha(x_1+x_2,y_1+y_2) = (\alpha(x_1+x_2),\alpha(y_1+y_2)) = (\alpha x_1,\alpha y_1) + (\alpha x_2,\alpha y_2) =  \alpha\vec{u} + \alpha\vec{v}$.
        \item $(\alpha + \beta)\vec{u} = ((\alpha+\beta)x_1,(\alpha+\beta)y_1) = (\alpha x_1 + \beta x_1, \alpha y_1 + \beta y_1) = (\alpha x_1,\alpha y_1) + (\beta x_1,\beta y_1) = \alpha\vec{u} + \beta\vec{u}$.
        \item $1\vec{u} = 1(x_1,y_1) = (x_1,y_1) = \vec{u}$.
        \item Primeiro temos:
        \[
            \alpha(\beta\vec{u}) = \alpha(\beta x_1,\beta y_1) = (\alpha(\beta x_1),\alpha(\beta y_1)) = ((\alpha\beta)x_1,(\alpha\beta)y_1) = (\alpha\beta)\vec{u}.
        \]
        Agora,
        \[
            (\alpha\beta)\vec{u} =     ((\alpha\beta)x_1,(\alpha\beta)y_1) = ((\beta\alpha)x_1,(\beta\alpha)y_1) = (\beta(\alpha x_1),\beta(\alpha y_1)) = \beta(\alpha\vec{u}).
        \]
        Logo $\alpha(\beta\vec{u}) = (\alpha\beta)\vec{u} = \beta(\alpha\vec{u})$.
    \end{enumerate}
    Portanto as propriedades são válidas.
\end{prova}

% \begin{proposicao}
%   Quaisquer que sejam o escalar $\alpha$ e o vetor $\vec{u}$, valem as seguintes igualdades:
%   \begin{enumerate}[label=({\roman*})]
%     \item $(-\alpha)\vec{u} = -(\alpha\vec{u})$;
%     \item $\alpha(-\vec{u}) = -(\alpha\vec{u})$;
%     \item $(-\alpha)(-\vec{u}) = \alpha\vec{u}$.
%   \end{enumerate}
% \end{proposicao}

Em alguns casos pode ser mais conveniente representar um vetor $\vec{u} = (x_1, y_1)$ de forma matricial. Para isso escrevemos as componentes de $\vec{u}$ como uma matriz de 1 coluna e duas linhas
\[
  \vec{u} = \begin{bmatrix}
    x_1\\y_1
  \end{bmatrix}.
\]

%Vimos que a norma de um vetor é definida como o comprimento de um segmento orientado que o represente. Com a representação de vetores por coordenadas podemos reescrever o conceito de norma da seguinte maneira:

\begin{exemplos}
  Seja $\vec{u} = (-2,3)$ um vetor. Então
  \begin{align*}
    \norm{\vec{u}} &= \sqrt{(-2)^2 + 3^2} = \sqrt{13}\\
    \vec{v} &= \dfrac{\vec{u}}{\norm{u}} = \dfrac{1}{\sqrt{13}}(-2,3) = \left(\dfrac{-2}{\sqrt{13}}, \dfrac{3}{\sqrt{13}}\right).
  \end{align*}
\end{exemplos}


% subsection operacoes_com_vetores (end)


% section vetores_plano (end)

\begin{proposicao}
  Sejam $\vec{u} = (x_1, y_1)$ um vetor em $\real^2$ e $\alpha \in \real$. Então:
  \begin{enumerate}[label=({\roman*})]
    \item Se $\vec{u} \ne \vec{0}$, então $\norm{\vec{u}} \ne 0$.
    \item $\norm{\vec{u}} = 0$ se, e somente se, $\vec{u} = \vec{0}$.
    \item $\norm{\alpha\vec{u}} = |\alpha|\norm{\vec{u}}$.
  \end{enumerate}
\end{proposicao}
\begin{prova}
  \begin{enumerate}[label=({\roman*})]
    \item  Se $\vec{u} \ne (0,0) = \vec{0}$, então $x_1 \ne 0$ ou $y_1 \ne 0$. Daí
    \[
      \norm{\vec{u}} = \sqrt{x_1^2 + y_1^2} > 0.
    \]
    \item $\norm{\vec{u}} = 0$ se, e somente se, $x_1^2 + y_1^2 = 0$, isto é, $x_1 = y_1 = 0$. Portanto $\vec{u} = \vec{0}$.
    \item $\norm{\alpha\vec{u}} = \norm{\alpha(x_1, y_1)} = \sqrt{(\alpha x_1)^2 + (\alpha y_1)^2} = \sqrt{\alpha^2(x_1^2 + y_1^2)} = |\alpha|\norm{\vec{u}}$.
  \end{enumerate}
\end{prova}

\begin{observacao}
    Dado um vetor não nulo $\vec{u}$ o vetor
    \[
          \vec{v} = \dfrac{\vec{u}}{\norm{\vec{u}}}
    \]
    é um vetor unitário na direção de $\vec{u}$ pois
    \[
          \norm{\vec{v}} = \norm{\dfrac{\vec{u}}{\norm{\vec{u}}}} = \dfrac{1}{\norm{\vec{u}}}\norm{\vec{u}} = 1.
    \]

\end{observacao}
\begin{definicao}
    Sejam $\vec{u}$ e $\vec{v}$ vetores não-nulos. Então $\vec{u} = \vec{v}$ se, e somente se, $\vec{u}$ e $\vec{v}$ têm normas iguais, são de mesma direção e de mesmo sentido.
\end{definicao}

\begin{proposicao}\label{vetores_paralelos}
  Dois vetores não nulos $\vec{u}$ e $\vec{v}$ são paralelos se, e somente se, existe um escalar $\lambda$ tal que $\vec{u} = \lambda\vec{v}$ (consequentemente, $\lambda \ne 0$ e assim $\vec{v} = \vec{u}/\lambda$).
\end{proposicao}
\begin{prova}
    Se $\vec{u} = \lambda\vec{v}$, então $\vec{u}\varparallel\vec{v}$.
    Agora, suponha que $\vec{u}$ e $\vec{v}$ são paralelos e de mesmo sentido, o caso em que têm sentidos contrários fica como exercício.

    Seja
    \[
        \lambda = \dfrac{\norm{\vec{u}}}{\norm{\vec{v}}}.
    \]
    Vamos mostrar que $\vec{u} = \lambda\vec{v}$. Inicialmente, sabemos que $\vec{v}$ e $\lambda\vec{v}$ são paralelos e têm o mesmo sentido pois $\lambda > 0$. Agora
    \[
        \norm{\lambda\vec{v}} = |\lambda|\norm{\vec{v}} = \dfrac{\norm{\vec{u}}}{\norm{\vec{v}}}\norm{\vec{v}} = \norm{\vec{u}}.
    \]
    Portanto $\vec{u}$ e $\lambda\vec{v}$ tem a mesma direção, o mesmo sentido e normas iguais, logo $\vec{u} = \lambda\vec{v}$, como queríamos.
\end{prova}

\subsection{ângulo entre vetores e produto interno} % (fold)
\label{sub:angulo_entre_vetores_e_produto_interno}
\begin{definicao}
  Sejam $\vec{u}$ e $\vec{v}$ vetores não nulos. O \textbf{ângulo} entre $\vec{u}$ e $\vec{v}$ é definido pelo ângulo $\theta$ determinado por $\vec{u}$ e $\vec{v}$ e que satisfaz $0 \le \theta \le \pi$, quando os vetores $\vec{u}$ e $\vec{v}$ são representados com a mesma origem.\index{Vetores!ângulo}
\end{definicao}

\begin{figure}[!h]
  \centering
  \caption{ângulo entre vetores}
  \includegraphics{angulo-vetores-plano-cartesiano.pdf}
\end{figure}


\begin{definicao}
  Quando o ângulo $\theta$ entre dois vetores $\vec{u}$ e $\vec{v}$ é reto, isto é, $\theta = \pi/2$, ou um deles é o vetor nulo, dizemos que os vetores $\vec{u}$ e $\vec{v}$ são \textbf{ortogonais} ou \textbf{perpendiculares} entre si. Denotamos tal fato, escrevendo $\vec{u} \perp \vec{v}$.\index{Vetores!Ortogonais}
\end{definicao}

Sejam $\vec{u}$ e $\vec{v}$ vetores não nulos e $\theta$ o ângulo entre eles.
\begin{figure}[!h]
  \centering
  \caption{Determinação de $\theta$}
  \includegraphics{determinacao-angulo-vetores-plano-cartesiano.pdf}
\end{figure}

Aplicando a Lei dos cossenos no triângulo abaixo
\begin{figure}[!h]
  \centering
  \caption{Lei dos cossenos}
  \includegraphics{lei-cossenos-vetores-plano-cartesiano.pdf}
\end{figure}
temos
\begin{align}\label{equacaonorma}
  \norm{\vec{u} - \vec{v}}^2 = \norm{\vec{u}}^2 + \norm{\vec{v}}^2 - 2\norm{\vec{u}}\norm{\vec{v}}\cos\theta.
\end{align}

% Assim conhecendo as coordenadas dos vetores $\vec{u}$ e $\vec{v}$ podemos determinar facilmente o ângulo entre eles. Além disso, note que se $\vec{u}$ e $\vec{v}$ são vetores ortogonais, então $\theta = \pi/2$ e então a equação \eqref{equacaonorma} torna-se
% \begin{align}\label{norma-vetores-ortogonais}
%   \norm{\vec{u} - \vec{v}} = \sqrt{\norm{\vec{u}}^2 + \norm{\vec{v}}^2}.
% \end{align}

% \begin{observacao}
%   A equação \eqref{norma-vetores-ortogonais} só é válida para vetores ortogonais entre si.
% \end{observacao}

\begin{definicao}\label{produtointerno}
  O \textbf{produto escalar} ou \textbf{produto interno} dos vetores $\vec{u}$ e $\vec{v}$, indicado por $\inner{u}{v}$, é o número real tal que:\index{Vetores!Produto Escalar}
  \begin{enumerate}
    \item Se $\vec{u}$ ou $\vec{v}$ é nulo, então $\inner{u}{v} = 0$.
    \item Se $\vec{u}$ e $\vec{v}$ não são nulos e $\theta$ é o ângulo entre $\vec{u}$ e $\vec{v}$, então
    \begin{align}\label{produto-interno}
      \inner{u}{v} = \norm{\vec{u}}\norm{\vec{v}}\cos\theta.
    \end{align}
  \end{enumerate}
\end{definicao}

\begin{proposicao}
  Sejam $\vec{u}$ e $\vec{v}$ vetores e $\theta$ o ângulo entre $\vec{u}$ e $\vec{v}$.
  \begin{enumerate}
    \item Se $\vec{u}$ e $\vec{v}$ não são nulos, então
    \[
      \cos\theta = \dfrac{\inner{u}{v}}{\norm{\vec{u}}\norm{\vec{v}}}.
    \]
    \item Qualquer que seja o vetor $\vec{u}$,
    \[
      \norm{\vec{u}} = \sqrt{\inner{u}{u}}.
    \]
    \item Quaisquer que sejam os vetores $\vec{u}$ e $\vec{v}$, $\vec{u}\perp\vec{v}$ se, e somente se, $\inner{u}{v} = 0$.
    \end{enumerate}
\end{proposicao}
\begin{prova}
  \begin{enumerate}
    \item Basta isolar $\cos\theta$ na Definição \ref{produto-interno}.
    \item Se $\vec{u} = \vec{0}$, então a igualdade é verdadeira. Se $\vec{u} \ne \vec{0}$, então o ângulo entre o vetor $\vec{u}$ e ele mesmo é $\theta = 0$, logo
    \[
      1 = \cos\theta = \dfrac{\inner{u}{v}}{\norm{\vec{u}}\norm{\vec{u}}}.
    \]
    Logo
    \[
      \norm{\vec{u}} = \sqrt{\inner{u}{u}}.
    \]
    \item Se $\vec{u}$ ou $\vec{v}$ é o vetor nulo, então $\vec{u}\perp\vec{v}$ e $\inner{u}{v} = 0$. Suponha $\vec{u} \ne \vec{0}$ e $\vec{v} \ne \vec{0}$. Assim, se $\vec{u} \perp\vec{v}$, então $\theta = \pi/2$ e daí $\inner{u}{v} = 0$. Agora, se $\inner{u}{v} = 0$, então
    \[
      0 = \inner{u}{v} = \norm{\vec{u}}\norm{\vec{v}}\cos\theta.
    \]
    Logo devemos ter $\cos\theta = 0$, isto é, $\theta = \pi/2$. Portanto $\vec{u}\perp\vec{v}$, como queríamos.
  \end{enumerate}
\end{prova}

Agora, como podemos encontrar o produto interno sem necessitar de determinar o ângulo $\theta$ entre dois vetores $\vec{u}$ e $\vec{v}$\,? Substituindo o produto interno na equação \eqref{equacaonorma} temos
\[
  \norm{\vec{u} - \vec{v}}^2 = \norm{\vec{u}}^2 + \norm{\vec{v}}^2 - 2\inner{u}{v}
\]
isto é,
\begin{align}\label{equacao-auxilar-norma}
  \inner{u}{v} = \dfrac{1}{2}(\norm{\vec{u}}^2 + \norm{\vec{v}}^2 - \norm{\vec{u} - \vec{v}}^2).
\end{align}
Se $\vec{u} = (x_1, y_1)$ e $\vec{v} = (x_2, y_2)$, então podemos reescrever \eqref{equacao-auxilar-norma} como
\begin{align*}
  \inner{u}{v} &= \dfrac{1}{2}(x_1^2 + y_1^2 + x_2^2 + y_2^2 - \norm{(x_1 - x_2, y_1 - y_2)}^2)\\
  &= \dfrac{1}{2}[x_1^2 + y_1^2 + x_2^2 + y_2^2 - (x_1 - x_2)^2 - (y_1 - y_2)^2]\\
  &= x_1x_2 + y_1y_2.
\end{align*}

Acabamos de demonstrar o seguinte teorema:
\begin{teorema}
  O produto interno de dois vetores $\vec{u} = (x_1, y_1)$ e $\vec{v} = (x_2, y_2)$ em $\real^2$ é dado por
  \[
    \inner{u}{v} = x_1x_2 + y_1y_2.
  \]
\end{teorema}

\begin{exemplos}
  \begin{enumerate}
    \item Sejam $\vec{u} = (1, 1)$ e $\vec{v} = (1,0)$. Determine o ângulo entre $\vec{u}$ e $\vec{v}$.
    \begin{solucao}
      Temos
      \begin{align*}
        \inner{u}{v} = \langle (1, 1), (1,0)\rangle = 2 + 0 = 1\\
        \norm{\vec{u}} = \sqrt{2},\ \norm{\vec{v}} = \sqrt{1}\\
        \cos\theta = \frac{\inner{u}{v}}{\norm{\vec{u}}\norm{\vec{v}}} = \dfrac{1}{\sqrt{2}}.
      \end{align*}
      Assim, $\theta = \pi/4$.
    \end{solucao}
    \item Encontre um vetor $\vec{u} = (x, y)$ ortogonal a $\vec{v} = (4, -1)$ e tal que $\inner{u}{w} = -1$, onde $\vec{w} = (1,1)$.
    \begin{solucao}
      Como $\inner{u}{v} = 0$ e $\inner{u}{w} = -1$ obtemos o sistema
      \[
        \begin{cases}
          4x - y = 0\\
          x + y = -1
        \end{cases}
      \]
      cuja solução é $x = -1/5$ e $y = -4/5$. Logo $\vec{u} = (-1/5, -4/5)$.
    \end{solucao}
  \end{enumerate}
\end{exemplos}

\begin{proposicao}\label{propriedades-produto-interno}
  Quaisquer que sejam os vetores $\vec{u} = (x_1, y_1)$, $\vec{v} = (x_2, y_2)$ e $\vec{w} = (x_3, y_3)$ e qualquer que seja o número real $\lambda$ temos
  \begin{enumerate}[label=({\roman*})]
    \item\label{linearidade-produto-interno} $\langle \vec{u}, \vec{v} + \vec{w}\rangle = \inner{u}{v} + \inner{u}{w}$
    \item $\langle \vec{u} + \vec{v}, \vec{w}\rangle = \inner{u}{w} + \inner{v}{w}$
    \item\label{linearidade2-produto-interno} $\langle (\lambda \vec{u}\,), \vec{v}\rangle = \langle \vec{u}, (\lambda \vec{v}\,)\rangle = \lambda\inner{u}{v}$
    \item $\inner{u}{v} = \inner{v}{u}$
    \item Se $\vec{u} \ne \vec{0}$, então $\inner{u}{u} > 0$.
  \end{enumerate}
\end{proposicao}
\begin{prova}
  \begin{enumerate}[label=({\roman*})]
    \item \begin{align*}
      \langle\vec{u}, (\vec{v} + \vec{w})\rangle &= \langle\vec{u}, (x_2 + x_3, y_2 + y_3)\rangle = x_1(x_2 + x_3) + y_1(y_2 + y_3) \\ &= (x_1x_2 + y_1y_2) + (x_1x_3 + y_1y_3) = \inner{u}{v} + \inner{u}{w}
    \end{align*}
    \item \begin{align*}
      \langle\vec{u}, (\lambda\vec{v})\rangle = \langle\vec{u}, (\lambda x_2, \lambda y_2)\rangle = x_1(\lambda x_2) + y_1(\lambda y_2) = (\lambda x_1)y_2 + (\lambda y_1)y_2 = \langle(\lambda\vec{u}),\vec{v}\rangle\\
      \langle\vec{u},(\lambda\vec{v})\rangle = \langle\vec{u}, (\lambda x_2, \lambda y_2)\rangle = x_1(\lambda x_2) + y_1(\lambda y_2) = \lambda (x_1y_2) + \lambda (y_1y_2) = \lambda\inner{u}{v}
    \end{align*}
    \item $\inner{u}{v} = x_1x_2 + y_1y_2 = x_2x_1 + y_2y_1 = \inner{v}{u}$
    \item Se $\vec{u} \ne \vec{0}$, então $x_1 \ne 0$ ou $y_1 \ne 0$. Daí $\inner{u}{u} = x_1^2 + y_1^2 > 0$, como queríamos.
  \end{enumerate}
\end{prova}

\begin{observacao}
  \begin{enumerate}
    \item As propriedades \ref{linearidade-produto-interno} e \ref{linearidade2-produto-interno} Proposição \ref{propriedades-produto-interno} podem ser estendidas para qualquer número de vetores e escalares:
    \[
      \langle\vec{u}, (\lambda_1\vec{v_1} + \lambda_2\vec{v_2} + \cdots + \lambda_n\vec{v_n} )\rangle = \lambda_1\inner{u}{v_1} + \lambda_2\inner{u}{v_2} + \cdots + \lambda_n\inner{u}{v_n}.
    \]
    \item Na igualdade $\inner{u}{v} = \inner{u}{w}$ não podemos concluir que $\vec{v} = \vec{w}$. Por exemplo, para $\vec{u} = (1, 0)$, $\vec{v} = (2, 1)$ e $\vec{w} = (2, -5)$ temos $\inner{u}{v} = \inner{u}{w}$ e no entanto $\vec{v} \ne \vec{w}$. Mas podemos concluir que $\vec{u}\perp(\vec{v} - \vec{w})$.
    \item De $\inner{u}{v} = 0$ não podemos concluir que $\vec{u} = \vec{0}$ ou $\vec{v} = \vec{0}$. Por exemplo, para $\vec{u} = (2, 1)$ e $\vec{v} = (-2, 4)$ temos $\inner{u}{v} = 0$ e no entanto $\vec{u} \ne \vec{0}$ e $\vec{v} \ne \vec{0}$.
  \end{enumerate}
\end{observacao}

\begin{exemplos}
  Calcule o ângulo $\alpha$ entre $\vec{u} + \vec{v}$ e $\vec{u} - \vec{v}$, sabendo que $\norm{\vec{u}} = \sqrt{5}$, $\norm{\vec{v}} = 1$ e que o ângulo $\theta$ entre $\vec{u}$ e $\vec{v}$ é $\pi/4$.
  \begin{solucao}
    Temos
    \begin{align*}
      \cos\alpha = \dfrac{\langle(\vec{u} + \vec{v}), (\vec{u} - \vec{v})\rangle}{\norm{\vec{u} + \vec{v}}\norm{\vec{u} - \vec{v}}} = \dfrac{\inner{u}{u} - \inner{u}{v} + \inner{v}{u} - \inner{v}{v}}{\norm{\vec{u} + \vec{v}}\norm{\vec{u} - \vec{v}}} = \dfrac{\norm{\vec{u}}^2 - \norm{\vec{v}}^2}{\norm{\vec{u} - \vec{v}}\norm{\vec{u} - \vec{v}}} = \dfrac{4}{\norm{\vec{u} + \vec{v}}\norm{\vec{u} - \vec{v}}}.
    \end{align*}
    Agora
    \begin{align*}
      \norm{\vec{u} + \vec{v}}^2 = \langle(\vec{u} + \vec{v}), (\vec{u} + \vec{v})\rangle = \inner{u}{v} + 2\inner{u}{v} + \inner{v}{v} = 6 + 2\inner{u}{v}\\
      \norm{\vec{u} - \vec{v}}^2 = \langle(\vec{u} - \vec{v}), (\vec{u} - \vec{v})\rangle = \inner{u}{v} - 2\inner{u}{v} + \inner{v}{v} = 6 - 2\inner{u}{v}.
    \end{align*}
    Por outro lado temos
    \begin{align*}
      \inner{u}{v} = \norm{\vec{u}}\norm{\vec{v}}\cos\theta = \sqrt{5}\cos(\pi/4) = \dfrac{\sqrt{10}}{2}
    \end{align*}
    e então
    \begin{align*}
      \norm{\vec{u} + \vec{v}} = \sqrt{6 + \sqrt{10}}\\
      \norm{\vec{u} - \vec{v}} = \sqrt{6 - \sqrt{10}}.
    \end{align*}
    Portanto
    \[
      \cos\theta = \dfrac{4}{\sqrt{26}}
    \]
    e então $\theta = \arccos\left(\dfrac{4}{\sqrt{26}}\right)$.
  \end{solucao}
\end{exemplos}

\begin{proposicao}\label{DesigualdadeTriangular}
  Sejam $\vec{u}$ e $\vec{v}$ vetores.
  \begin{enumerate}
    \item $|\inner{u}{v}| \le \norm{\vec{u}}\norm{\vec{v}}$ [Desigualdade de Scharwz]\index{Desigualdade de Scharwz}
    \item $\norm{\vec{u} + \vec{v}} \le \norm{\vec{u}} + \norm{\vec{v}}$ [Desigualdade Triangular]\index{Desigualdade Triangular}
  \end{enumerate}
\end{proposicao}
\begin{prova}
  \begin{enumerate}
    \item Se $\vec{u} = \vec{0}$ ou $\vec{v} = \vec{0}$ a desigualdade é verdadeira. Suponha então que $\vec{u} \ne \vec{0}$ e $\vec{v} \ne \vec{0}$. Seja $\theta$ o ângulo entre $\vec{u}$ e $\vec{v}$. Temos
    \[
      0 \le |\cos\theta| \le 1
    \]
    e multiplicando essa inequação por $\norm{\vec{u}}\norm{\vec{v}}$ obtemos
    \[
      0 \le \norm{\vec{u}}\norm{\vec{v}}\cos\theta \le \norm{\vec{u}}\norm{\vec{v}},
    \]
    isto é,
    \[
      |\inner{u}{v}| \le \norm{\vec{u}}\norm{\vec{v}}
    \]
    como queríamos.
    \item Se $\vec{u} = \vec{0}$ ou $\vec{v} = \vec{0}$, então a desigualdade é verdadeira. Suponha então que $\vec{u} \ne \vec{0}$ e $\vec{v} \ne \vec{0}$. Temos
    \begin{align*}
      \norm{\vec{u} + \vec{v}}^2 &= \langle\vec{u} + \vec{v}, \vec{u} + \vec{v}\rangle = \inner{u}{v} + 2\inner{u}{v} + \inner{v}{v} = \norm{\vec{u}}^2 + 2\inner{u}{v} + \norm{\vec{v}}^2 \\ &\le \norm{\vec{u}}^2 + 2\mid\inner{u}{v}\mid + \norm{\vec{v}}^2.
    \end{align*}
    Agora, usando o item (a) temos
    \begin{align*}
      \norm{\vec{u} + \vec{v}}^2 \le \norm{\vec{u}} + 2\mid\inner{u}{v}\mid \norm{\vec{v}}^2 \le \norm{\vec{u}} + 2\norm{\vec{u}}\norm{\vec{v}} + \norm{\vec{v}}^2 = (\norm{\vec{u}} + \norm{\vec{v}})^2.
    \end{align*}
    Portanto
    \[
      \norm{\vec{u} + \vec{v}} \le \norm{\vec{u}} + \norm{\vec{v}}.
    \]
  \end{enumerate}
\end{prova}

% subsection angulo_entre_vetores_e_produto_interno (end)
\subsection{Projeção de Vetores} % (fold)
\label{sub:projecao_de_vetores}
Sejam $\vec{u} = (x_1, y_1)$ e $\vec{v} = (x_2, y_2)$ vetores não nulos. Suponha que o ângulo $\theta$ entre $\vec{u}$ e $\vec{v}$ é tal que $0 < \theta < \pi/2$.
\begin{figure}[!h]
  \centering
  \caption{Projeção ortogonal}
  \includegraphics{projecao-vetores-plano-cartesiano.pdf}
\end{figure}

O vetor $\vec{OP}$ é paralelo ao vetor $\vec{u}$ e o vetor $\vec{PA}$ é ortogonal ao vetor $\vec{u}$.

\begin{definicao}\label{projecao_ortogonal}
  Seja $\vec{u}$ um vetor não nulo. Dado qualquer vetor $\vec{v}$, o vetor $\vec{OP}$ é chamado de \textbf{projeção ortogonal} de $\vec{v}$ sobre $\vec{u}$, e é indicado por $\proj_{\vec{u}}\vec{v}$ se satisfaz as seguintes condições:\index{Vetores!Projeção ortogonal}
  \begin{enumerate}
    \item $\vec{OP} \varparallel\vec{u}$
    \item $(\vec{v} - \vec{OP}) \perp \vec{u}$.
  \end{enumerate}
\end{definicao}

Podemos também escrever
\begin{enumerate}
  \item $\proj_{\vec{u}}\vec{v} \varparallel\vec{u}$
    \item $(\vec{v} - \proj_{\vec{u}}\vec{v}) \perp \vec{u}$.
\end{enumerate}

No caso acima temos
\[
  \cos\theta = \dfrac{\norm{\vec{OP}}}{\norm{\vec{v}}}
\]
e então $\norm{\vec{OP}} = \norm{\vec{v}}\cos\theta$. Mas
\[
  \cos\theta = \dfrac{\inner{u}{v}}{\norm{\vec{u}}\norm{\vec{v}}}
\]
e assim podemos escrever
\[
  \norm{\vec{OP}} = \dfrac{\inner{u}{v}}{\norm{\vec{u}}}.
\]
Além disso, como $\vec{OP}$ é paralelo à $\vec{u}$ e de mesmo sentido, devemos ter $\vec{OP} = \lambda\vec{u}$ com $\lambda > 0$. Daí
$\lambda = \norm{\vec{OP}}/\norm{\vec{u}}$ e então
\[
  \vec{OP} = \dfrac{\norm{OP}}{\norm{\vec{u}}^2}\vec{u}.
\]

Esta construção só é válida no caso em que o ângulo $\theta$ entre $\vec{u}$ e $\vec{v}$ satisfaz $0 < \theta < \pi/2$. Nos casos
\begin{figure}
  \centering
  \caption{Vetor projeção}
  \includegraphics{vetores-projecao-plano-cartesiano.pdf}
\end{figure}
não podemos usar a construção envolvendo o ângulo $\theta$ entre os vetores.
\begin{proposicao}
  Seja $\vec{u}$ um vetor não nulo. Qualquer que seja o vetor $\vec{v}$, existe e é única a projeção ortogonal de $\vec{v}$ sobre $\vec{u}$. Sua expressão em termos de $\vec{u}$ e $\vec{v}$ é dada por
  \[
    \proj_{\vec{u}}\vec{v} = \dfrac{\inner{v}{u}}{\norm{\vec{u}}^2}\vec{u}
  \]
  e sua norma é dada por
  \[
    \norm{\proj_{\vec{u}}\vec{v}} = \dfrac{\mid\inner{v}{u}\mid}{\norm{\vec{u}}}.
  \]
\end{proposicao}
\begin{prova}
  Seja $\vec{p} = \proj_{\vec{u}}\vec{v}$. Da Definição \ref{projecao_ortogonal}, dizer que $\vec{p}$ é a projeção ortogonal de $\vec{v}$ sobre $\vec{u}$ significa que $\vec{p}$ é um vetor tal que $\vec{p}\varparallel\vec{u}$ e $(\vec{v} - \vec{p})\perp\vec{u}$.

  Como $\vec{p}\varparallel\vec{u}$, então pela Proposição \ref{vetores_paralelos} existe um escalar $\lambda$ tal que $\vec{p} = \lambda\vec{u}$. Assim qualquer vetor do tipo $\lambda\vec{u}$ é um candidato para ser a projeção ortogonal. Mas, para que $\vec{p}$ seja a projeção ortogonal, ele deve satisfazer $(\vec{v} - \vec{p})\perp\vec{u}$. Logo
  \begin{align*}
    (\vec{v} - \vec{p})\perp\vec{u} \Leftrightarrow \langle(\vec{v} - \vec{p}), \vec{u}\rangle = 0 \Leftrightarrow \inner{v}{u} - \lambda\inner{u}{v} = 0 \Leftrightarrow \lambda = \dfrac{\inner{v}{u}}{\norm{\vec{u}}^2}.
  \end{align*}
  Assim existe um único $\lambda$ que satisfaz a condição $(\vec{v} - \vec{p})\perp\vec{u}$, e daí o vetor $\vec{p}$ é único. Portanto
  \[
    \proj_{\vec{u}}\vec{v} = \vec{p} = \dfrac{\inner{v}{u}}{\norm{\vec{u}}^2}\vec{u}.
  \]
  Além disso
  \[
    \norm{\proj_{\vec{u}}\vec{v}} = \norm{\dfrac{\inner{v}{u}}{\norm{\vec{u}}^2}\vec{u}} = \dfrac{\mid\inner{v}{u}\mid}{\norm{\vec{u}}}
  \]
  como queríamos.
\end{prova}

\begin{observacao}
  O escalar $\lambda$ encontrado na proposição anterior pode ser negativo ou zero. Ele será negativo quando $\pi/2 \le \theta \le \pi$ e será zero quando $\vec{u}\perp\vec{v}$.
  \begin{figure}[!h]
    \centering
    \caption{Possibilidades para $\lambda$}
    \includegraphics{vetores-projecao-escalar-plano-cartesiano.pdf}
  \end{figure}
\end{observacao}

\begin{exemplos}\label{exemplosprojecao}
  \begin{enumerate}
    \item Encontre a projeção ortogonal do vetor $\vec{v} = (4, -1)$ sobre o vetor $\vec{u} = (2, 1)$.
    \begin{solucao}
       Temos
       \[
          \proj_{\vec{u}}\vec{v} = \dfrac{\inner{v}{u}}{\norm{\vec{u}}^2}\vec{u}.
       \]
       Agora
       \begin{align*}
         \inner{v}{u} = 7\\
         \norm{\vec{u}}^2 = 5
       \end{align*}
       e daí
       \[
          \proj_{\vec{u}}\vec{v} = \dfrac{7}{5}(2, 1) = \left(\dfrac{14}{5}, \dfrac{7}{5}\right).
       \]
       Por outro lado,
       \[
          \proj_{\vec{v}}\vec{u} = \dfrac{\inner{u}{v}}{\norm{\vec{v}}^2}\vec{v}.
       \]
       Mas
       \begin{align*}
         \inner{u}{v} = \inner{v}{u} = 7\\
         \norm{\vec{v}}^2 = 17
       \end{align*}
       e daí
       \[
          \proj_{\vec{v}}\vec{u} = \dfrac{7}{17}(4, -1) = \left(\dfrac{28}{17}, \dfrac{-7}{17}\right).
       \]
     \end{solucao}
     \item Considere um triângulo $ABC$ onde $A(3,3)$, $B(0,1)$ e $C(1,6)$.
     \begin{enumerate}\label{exemploprojecao2}
       \item Verifique que $ABC$ é um triângulo retângulo em $A$.
       \item Calcule a projeção do cateto $BA$ sobre a hipotenusa $BC$.
       \item Determine o pé da altura do triângulo relativo ao vértice $A$.
     \end{enumerate}
     \begin{solucao}
       Considere o triângulo abaixo
       \begin{figure}[!h]
        \centering
        \caption{Exemplo \ref{exemplosprojecao}, exercício \ref{exemploprojecao2}}
        \includegraphics{exemplo-vetores-projecao-plano-cartesiano.pdf}
       \end{figure}
       \begin{enumerate}
         \item Basta verificar que $\vec{AB}\perp\vec{AC}$. Assim como $\vec{AB} = (-3,-2)$ e $\vec{AC} = (-2,3)$ temos
         \[
            \inner{AB}{AC} = (-3)(-2) + (-2)3 = 0.
         \]
         Logo $ABC$ é um triângulo retângulo.
         \item Temos $\vec{BA} = (3,2)$ e $\vec{BC} = (1,5)$. Daí
         \[
            \proj_{\vec{BC}}\vec{BA} = \dfrac{\inner{BC}{BA}}{\norm{\vec{BC}}^2}\vec{BC} = \dfrac{\langle (3,2), (1,5)\rangle}{\norm{(1,5)}^2}(1,5) = \dfrac{13}{26}(1,5).
         \]
         Assim $\proj_{\vec{BC}}\vec{BA} = \left(\dfrac{1}{2},\dfrac{5}{2}\right)$.
         \item Seja $P(x_1, y_1)$ o pé da altura do triângulo $ABC$ relativo ao vértice $A$. Então
         \begin{align*}
           \vec{BP} = \proj_{\vec{BC}}\vec{BA}\\
           (x_1, y_1 - 1) = \left(\dfrac{1}{2}, \dfrac{5}{2}\right).
         \end{align*}
         Logo $P\left(\dfrac{1}{2},\dfrac{7}{2}\right)$.
       \end{enumerate}
     \end{solucao}
  \end{enumerate}
\end{exemplos}
% subsection projecao_de_vetores (end)
% chapter vetores (end)
